\section{Introduction}\label{sec:intro}

Multi-object tracking (MOT) finds every object of interest in a video and
keeps the same identity for each object over time. It is a core part of
surveillance, driving, and airborne systems
\cite{zhu2022visdrone,wu2022uavsurvey,lu2024uavdetection}, and many of these systems
must produce results with low delay. Most practical trackers follow the
tracking-by-detection design. A detector scores candidate boxes in each frame,
and the tracker links the boxes whose scores pass its thresholds into
trajectories \cite{bewley2016sort,zhang2022bytetrack,cao2023ocsort}.

In deployed systems the detector and the tracker are rarely developed
together. A detector is replaced by a faster model, retrained for a new
sensor, or quantized for embedded hardware. The tracker, however, keeps the
thresholds that its authors tuned for another detector. Detector confidence is
not calibrated across models and conditions \cite{kuppers2020detcalib}, so the
new detector usually has a different score scale, and nothing in the pipeline
notices the change.

The effect can be large. In the experiments of this paper, halving the scores
of a published detector drives two well-known trackers to 0 higher order
tracking accuracy (HOTA) \cite{luiten2021hota}, because no detection reaches
their fixed thresholds (Sect.~\ref{sec:shift}). Retuning
every tracker for every detector needs labelled data from the deployment
domain, which an operational system usually does not have.

Existing methods only partly address this problem. Adaptive-threshold
trackers adjust their own thresholds from the detections of each frame, but
each method is built into one specific tracker
\cite{ma2023adaptivebytetrack,shim2024adaptrack,act2026adaptive}. Confidence
calibration methods correct the detector scores, but they need labelled data
\cite{kuppers2020detcalib}. Adaptive inference methods learn or profile a
policy for one model \cite{han2022dynamic,jiang2018chameleon}. What is missing
is a training-free component that works outside the tracker, calibrates itself
from the detector's own score stream, and leaves the tracker unchanged when
the tracker already fits that stream.

This paper presents such a component: \scl{}, a self-calibrating control layer
between a frozen detector and a frozen tracker. It reads the stream of detector scores
and the operating point that the tracker declares, and it decides online
whether the tracker needs help. \scl{} is described in
Sect.~\ref{sec:method}, and Sect.~\ref{sec:development} explains how it was
developed from an earlier design that always intervened.

The main contributions are:
\begin{enumerate}
\item \scl{}, a training-free control layer that calibrates itself from the detector
      score stream and changes the tracker's input only when the stream is
      noisy (Sect.~\ref{sec:method}).
\item A host contract of four numbers that lets one frozen configuration work
      with single-stage, two-stage, coordinate-only, and two-view trackers
      without modifying them (Sect.~\ref{sec:contract}).
\item A development path from always-on banding (\aob{}), a prior design that
      always intervened, in which
      each step is compared with that baseline (Sect.~\ref{sec:development}).
\item An evaluation of one frozen configuration on nine published trackers,
      four detection sources, and three datasets, with development,
      predeclared, and post-freeze roles, paired bootstrap intervals, and a
      latency measurement on a CPU (Sects.~\ref{sec:setup}
      and~\ref{sec:results}).
\end{enumerate}

The rest of the paper is organized as follows. Section~\ref{sec:related}
reviews related work. Section~\ref{sec:method} describes \scl{}, and
Section~\ref{sec:development} explains how it was developed.
Section~\ref{sec:setup} gives the experimental setup,
Section~\ref{sec:results} the results, and Section~\ref{sec:discussion} a
discussion. Section~\ref{sec:conclusion} concludes the paper.
